\documentclass[a4paper,11pt]{article}

% ---------- Кодировка и язык ----------
\usepackage[utf8]{inputenc}
\usepackage[T2A]{fontenc}
\usepackage[english,russian]{babel}

% ---------- Геометрия страницы ----------
\usepackage[margin=2.5cm]{geometry}

% ---------- Математика ----------
\usepackage{amsmath,amssymb,amsthm,mathtools}
\usepackage{bm}          % жирные символы в формулах: \bm{x}

% ---------- Картинки ----------
\usepackage{graphicx}
\graphicspath{{figures/}}   % картинки лежат в latex/figures/
\usepackage{float}          % позиция [H] для картинок
\usepackage{subcaption}     % несколько картинок рядом

% ---------- Оформление ----------
\usepackage{enumitem}
\usepackage{booktabs}       % красивые таблицы
\usepackage{xcolor}
\usepackage[colorlinks=true,linkcolor=blue,urlcolor=blue,citecolor=blue]{hyperref}
\usepackage{cleveref}       % \cref{...} для умных ссылок
\newcommand{\lean}[1]{\mbox{\path{#1}}} % имена из Lean-кода, без переносов внутри имени (пакет url)

% ---------- Теоремы и окружения ----------
\theoremstyle{plain}
\newtheorem{theorem}{Теорема}[section]
\newtheorem{lemma}[theorem]{Лемма}
\newtheorem{proposition}[theorem]{Утверждение}
\newtheorem{corollary}[theorem]{Следствие}

\theoremstyle{definition}
\newtheorem{definition}[theorem]{Определение}
\newtheorem{example}[theorem]{Пример}

\theoremstyle{remark}
\newtheorem*{remark}{Замечание}

% ---------- Полезные макросы ----------
\newcommand{\R}{\mathbb{R}}
\newcommand{\Z}{\mathbb{Z}}
\newcommand{\N}{\mathbb{N}}
\newcommand{\Cx}{\mathbb{C}}   % \C занят babel-russian
\newcommand{\Tau}{\mathcal{T}} % тайлинг
\newcommand{\norm}[1]{\left\lVert #1 \right\rVert}
\newcommand{\abs}[1]{\left\lvert #1 \right\rvert}
\newcommand{\set}[1]{\left\{ #1 \right\}}
\newcommand{\ip}[2]{\langle #1, #2\rangle} % скалярное произведение
\newcommand{\Rpp}[1][2]{\R^{#1}_{++}} % \Rpp = R^2_{++}, \Rpp[n] = R^n_{++}
\newcommand{\Rp}[1][2]{\R^{#1}_{+}}   % \Rp = R^2_{+}, \Rp[n] = R^n_{+}
\DeclareMathOperator{\argmin}{arg\,min}
\DeclareMathOperator{\argmax}{arg\,max}

% Заметка на полях / TODO
\newcommand{\todo}[1]{\textcolor{red}{\textbf{TODO:} #1}}

\begin{document}
\section{Разрешимость}
\begin{definition}[Неоклассические Функции]
    Пусть $n\in\mathbb{N}$
    Будем называть функцию $f:\Rp[n]\to\R_{+}$ \textit{неоклассической}, если она \textit{непрерывна}, \textit{вогнута} и однородна степени~1.\\
    Будем обозначать класс неоклассических функций $n$ переменных как $\Phi_n$.
\end{definition}

\begin{definition}[CES Функции]
    Пусть $n\in\mathbb{N}$, $J\subset[n]$, $\epsilon\in[-1,0)\cup(0,+\infty)$. Тогда отображение $A:J\to\R_{++}$ задает параметрическое семейство неоклассических функций
    \[H(X) := \Bigl(\sum_{j\in J}\Bigl(\frac{X_{j}}{A_{j}}\Bigr)^{-\epsilon}\Bigr)^{-1/\epsilon},\]
    называемых функциями \textit{постоянной эластичности замещения} $\sigma := \frac{1}{1+\epsilon}$.
    При $\epsilon\to 0$ семейство сходится к функции Кобба--Дугласа ($\sigma=1$), при $\epsilon\to+\infty$ --- к леонтьевской $\min_{j\in J} X_j/A_j$ ($\sigma=0$), при $\epsilon=-1$ --- к линейной $\sum_{j\in J} X_j/A_j$ ($\sigma=\infty$). При $\epsilon<-1$ функция не вогнута и потому не является неоклассической. Обозначим класс CES функций $\Psi_n\subset\Phi_n$.
\end{definition}

\begin{definition}[Неоклассическая разрешимость]\label{def:neo-solvability}
    Пусть $\{\hat p_t \in \Rpp\}_{t\in[T]}$ --- некоторый конечный набор цен.
    Будем говорить, что вектор выпусков $y\in\Rpp[T]$ \textit{неоклассически разрешим} при этих ценах, если
    существует такая неоклассическая функция $h\in\Phi_{2}$ и такие окрестности
    $\{U_t\ni\hat p_t\}_{t\in[T]}$, что для каждого набора $\{p_t\in U_t\}_{t\in[T]}$
    существует полная абсолютно непрерывная мера $\mu$ на $\Rp$, что $\forall t\in[T]\Rightarrow\mu(A_t)=y_t$, где ${A_t=\{x\in\Rpp\mid h(p_t\circ x)<1\}}.$
\end{definition}

\begin{definition}[CES-разрешимость]
    Будем говорить, что вектор выпусков $y\in\Rpp[T]$ \textit{CES-разрешим} при ценах $\{\hat p_t\}_{t\in[T]}$, если в предыдущем определении функцию $h$ можно выбрать из класса $\Psi_{2}$. Поскольку $\Psi_{2}\subset\Phi_{2}$, из CES-разрешимости следует неоклассическая разрешимость.
\end{definition}

\begin{proposition}
    $\exists q\in conv\{p_{1},p_{3}\}:q\leq p_{2}\Rightarrow A_{2}\subset A_{1}\cup A_{3}$.
\end{proposition}

\begin{example}
    При ценах $\{\hat{p}_{1} = (4, 1), \hat{p}_{2} = (3, 3), \hat{p}_{3} = (1, 4)\}$ вектор выпусков $y=(1, 5, 3)$ не является неоклассически разрешимым (а значит и CES-разрешимым).
    Более того, для случая CES функций, критическое значение параметра $\epsilon_{crit}\approx -2$, а значит ромбический тайлинг будет иметь вид $\lambda$-гекса независимо от эластичности замещения. Вектору выпусков $y=(1, 5, 3)$ соответствует змейка $(0,\xi_{2},\xi_{2}+\xi_{3},\xi_{1}+\xi_{2}+\xi_{3})$ которая лежит внутри зоногона $Z_{3}$ но не содержится в $\lambda$-гексе. 
\end{example}

\begin{example}
    При ценах $\hat p_1=(1,9)$, $\hat p_2=(4,4)$, $\hat p_3=(9,1)$, $\hat p_4=(10,100000)$, $\hat p_5=(20000,20000)$, $\hat p_6=(100000,10)$
    ромбический тайлинг любой CES функции состоит из двух гексов, склеенных в вершине со спектром $\{1,2,3\}$, каждый из которых имеет вид $\lambda$ или $\gamma$. Нижний гекс флипается при $\epsilon\approx-0.50$, а верхний --- при $\epsilon\approx-0.42$, поэтому при проходе $\epsilon$ по области допустимых значений возможны всего три спектра $Sp_{\lambda\lambda}$, $Sp_{\gamma\lambda}$, $Sp_{\gamma\gamma}$, а четвёртый спектр $Sp_{\lambda\gamma}$ невозможен. Вектор выпусков $y=(5,9,6,3,1,4)$ лежит в конусе $C(Sp_{\lambda\gamma})$ и не лежит ни в одном из конусов $C(Sp_{\lambda\lambda})$, $C(Sp_{\gamma\lambda})$, $C(Sp_{\gamma\gamma})$, а значит, не является CES-разрешимым.
    В то же время для неоклассической функции ${h(x):=(x^{1}x^{2})^{0.15}(x^{1}+x^{2})^{0.7}}$ линии уровня $h(\hat{p}_{t}\circ x)=1$ образуют строгую проволочную диаграмму с невозможным для CES функций спектром $Sp_{\lambda\gamma}$, значит $y$ неоклассически разрешим.
\end{example}

\newpage
\section{Конечность Пересечений}

\begin{proposition} Функция $h(x):=\sqrt{x^{1}x^{2}}\exp(\frac{1}{4}\sin\ln\frac{x^{1}}{x^{2}})$ является неоклассической.
\end{proposition}

\begin{proof}
    Сначала покажем, что эта функция непрерывна на $\Rp$. Непрерывность на $\Rpp$ выполнена в силу того, что $h$ представима как композиция непрерывных функций. Для того, чтобы показать непрерывность на $\partial\Rp$ заметим, что  для каждого $x\in\Rpp$ значение функции $h(x)\in[e^{\frac{-1}{4}}\sqrt{x^{1}x^{2}}, e^{\frac{+1}{4}}\sqrt{x^{1}x^{2}}]$ а значит также как и в случае функций Кобба-Дугласса, наша функция продолжается по непрерывности на $\partial\Rp$ тождественным нулем.

    Теперь покажем, что эта функция вонута. Введем $u:=\ln\frac{x^{1}}{x^{2}}$, тогда гессиан имеет вид:
    \[
        \nabla^{2}h(x)
        = -\,\frac{h(x)\,\bigl(1+\sin u\bigr)\bigl(3+\sin u\bigr)}{16\,(x^{1})^{2}(x^{2})^{2}}
        \begin{pmatrix}
            (x^{2})^{2} & -x^{1}x^{2}\\[2pt]
            -x^{1}x^{2} & (x^{1})^{2}
        \end{pmatrix}.
    \]
    Множитель перед матрицей неположителен, а сама матрица имеет собственные числа $0$ и $(x^{1})^{2}+(x^{2})^{2}$, а значит сам гессиан отрицательно полуопределен и функция $h$ вогнута.
\end{proof}

\begin{example}
    Пусть $h(x)=\sqrt{x^{1}x^{2}}\exp(\frac{1}{4}\sin\ln\frac{x^{1}}{x^{2}})$ и $p\in\Rpp$, тогда если выполнено неравенство $|\ln p^{1} + \ln p^{2}|<|\sin\frac{\ln p^{1} - \ln p^{2}}{2}|$ то кривые $h(x)=1$ и $h(p\circ x)=1$ имеют счетное пересечение.
\end{example}

\begin{proof}
    Введем замену координат $\theta:= \ln x^{1} + \ln x^{2}$ и $\eta := \ln x^{1} - \ln x^{2}$, в этих координатах кривая $h(x)=1$ имеет простой параметрический вид $\theta =  -\frac{\sin\eta}{2}$ и определена для любого $\eta$.

    Помимо равенства $h(x)=1$ для точки пересечения также требуется $h(x)=h(p\circ x)$.
    Полагая, что $h(x)>0$ в координатах $(\theta,\eta)$ это уравнение представимо в следующем виде:
    \[
    2\ln\frac{h(p\circ x)}{h(x)}
    = (\ln p^{1} + \ln p^{2})
    + \sin\frac{\ln p^{1} - \ln p^{2}}{2}\,\cos\left(\eta + \frac{\ln p^{1} - \ln p^{2}}{2}\right) = 0.
    \]
    Это уравнение зависит от $\theta$ и имеет $0$ корней при $|\ln p^{1} + \ln p^{2}|>|\sin\frac{\ln p^{1} - \ln p^{2}}{2}|$, континуум решений в вырожденных случаях $p = (\exp(\pi z), \exp(-\pi z))$ для $z\in\mathbb{Z}$ и счетное кол-во в остальных случаях в том числе при $|\ln p^{1} + \ln p^{2}|<|\sin\frac{\ln p^{1} - \ln p^{2}}{2}|$. Для любого $\eta$ существует единственная $\theta$, которая задает точку на $h(x)=1$, что дает счетное кол-во пересечений. Более того, множество таких параметров $p$ открыто, а значит счетное кол-во пересечений достигается в \textit{общем положении}.
\end{proof}

\begin{definition}[Строго Положительная Неоклассическая Функция] Будем называть неоклассическую функцию $h\in\Phi_{n}$ \textit{строго положительной} если $\forall p\in\Rp[n]\setminus\set{0}\Rightarrow h(p)>0$. Класс строго положительных функций $n$ переменных будем обозначать как $\hat{\Phi}_{n}$.
\end{definition}

\begin{example} Пусть $\set{\xi_{j}\in\Rpp}_{j\in J}$ - конечное множество векторов, тогда функция\\
    $h(p):=\min\limits_{j\in J}\langle\xi_{j},p\rangle$
    является строго положительной неоклассической функцией.
\end{example}

\begin{definition}[Конус Мордуховича~\cite{Mordukhovich2006}]
Пусть $A\subset\R^n$ замкнуто, $\bar x\in A$. \\
\emph{регулярный нормальный конус Фреше}:
\[\widehat N_A(\bar x):=\Bigl\{v\in\R^n\Bigm|
\limsup_{x\to\bar x,x\in A,x\ne\bar x}
\frac{\ip{v}{x-\bar x}}{\abs{x-\bar x}}\le 0\Bigr\},\]
\emph{предельный нормальный конус Мордуховича}:
\[N_A(\bar x):=\bigl\{v\in\R^n\bigm|\exists x_k\to\bar x,\ x_k\in A,\
v_k\to v,\ v_k\in\widehat N_A(x_k)\bigr\}.\]
Для выпуклого $A$ оба конуса совпадают с классическим определением нормального конуса.
\end{definition}

\begin{definition}[Трансверсальность~\cite{Kruger2006,Ioffe2017}]
Пусть $A,B\subset\R^n$ замкнуты, $x_{0}\in A\cap B$. Пара $\{A,B\}$
\emph{трансверсальна} в точке $x_{0}$, если
$N_A(x_{0})\cap\bigl(-N_B(x_{0})\bigr)=\{0\}.$
Набор $\{A_1,\dots,A_k\}$ трансверсален в $x_{0}\in\bigcap_i A_i$, если из
$v_1+\dots+v_k=0$, $v_i\in N_{A_i}(x_{0})$, следует $v_1=\dots=v_k=0$.
\end{definition}

\begin{proposition}\label{prop:transversal} Пусть $h,g\in\hat{\Phi}_{2},x_{0}\in\Rpp:h(x_{0})=g(x_{0})=1$. Линии уровня $h(x)=1$ и $g(x)=1$ пересекаются в $x_{0}$ трансверсально, если и только если $\partial h(x_{0})\cap\partial g(x_{0}) = \varnothing.$
Более того, точка $x_{0}\in\partial\Rpp$ не может быть трансверсальным пересечением этих кривых.
\end{proposition}

\begin{remark}
В дальнейшем эквивалентное условие из утверждения~\ref{prop:transversal} мы используем в качестве определения: точку $x_0$ пересечения линий уровня $h(x)=1$ и $g(x)=1$ называем \emph{трансверсальной}, если $x_0\in\Rpp$ и $\partial h(x_0)\cap\partial g(x_0)=\varnothing$.
\end{remark}

\begin{definition}[Общее положение относительно $h$]\label{def:general-position}
    Пусть $h\in\hat\Phi_{2}$ и $P=\{p_t\in\Rpp\}_{t\in[T]}$ --- набор цен. Для $s,t\in[T]$ обозначим через $I_{st}(P):=\set{x\in\Rp\mid h(p_s\circ x)=1,\ h(p_t\circ x)=1}$ множество точек пересечения линий уровня $h(p_s\circ x)=1$ и $h(p_t\circ x)=1$.
    Будем говорить, что пара $(h,P)$ находится \textit{в общем положении}, если выполнены следующие пункты:
    \begin{enumerate}
        \item для любой пары различных индесов $s\neq t$ линии уровня $h(p_{s}\circ x)=1$ и $h(p_{t}\circ x)=1$ имеют конечное кол-во пересечений, причем все эти пересечения трансверсальны.
        \item ни для какой тройки различных индексов $r\neq s\neq t\neq r$ линии уровня 
        \\ $h(p_{r}\circ x)=1, h(p_{s}\circ x)=1$ и $h(p_{t}\circ x)=1$ не пересекаются в одной точке.
    \end{enumerate}
\end{definition}

\begin{lemma}\label{thm:finite} Пусть $h\in\hat{\Phi}_{2}$, и $G_h\subset\Rpp$ --- множество цен $p$, при которых кривые $h(x)=1$ и $h(p\circ x) = 1$ имеют конечное число пересечений, причём все эти пересечения происходят трансверсально. Тогда $G_h$ открыто и всюду плотно в $\Rpp$.
\end{lemma}

\begin{proposition}[Устойчивость общего положения]\label{prop:gp-stable}
    Пусть $h\in\hat\Phi_{2}, P\in\Rpp[2\times T]$ и пара $(h,P)$ находится в общем положении.
    Тогда существует такая окрестность $V\subset\Rpp[2\times T]$ набора $P$, что для каждого набора
    цен $Q\in V$ пара $(h,Q)$ также в общем положении.
\end{proposition}

\begin{proof}
Пункт~1 определения~\ref{def:general-position} сохраняется вблизи $P$ по лемме~\ref{thm:finite}: для пары $(s,t)$ он означает $p_t/p_s\in G_h$, а $G_h$ открыто.
Пусть пункт~2 нарушается для наборов $Q_n\to P$, каждый раз для одних и тех же индексов $r,s,t$. Напомним, что $I_{st}(Q)$ --- множество точек пересечения линий уровня $h(q_s\circ x)=1$ и $h(q_t\circ x)=1$ в $\Rp$ для набора цен $Q=\{q_t\}_{t\in[T]}$. При ценах, близких к $P$, все точки пересечения лежат в одном ограниченном множестве: из $1=h(q_s\circ x)\ge h(1,0)q_s^1x^1+h(0,1)q_s^2x^2$ следует $x^1\le1/(h(1,0)q_s^1)$ и $x^2\le1/(h(0,1)q_s^2)$, а $q_s^1,q_s^2$ отделены от нуля. Поэтому из тройных точек $x_n\in I_{rs}(Q_n)\cap I_{rt}(Q_n)$ можно выбрать сходящуюся подпоследовательность, и её предел лежит в $I_{rs}(P)\cap I_{rt}(P)$ по непрерывности $h$, то есть является тройной точкой для $P$, что противоречит общему положению $(h,P)$.
\end{proof}

\begin{proposition}[Плотность общего положения]\label{prop:gp-dense}
    Пусть $h\in\hat{\Phi}_{2}$. Тогда в любом непустом открытом множестве $U\subset\Rpp[2\times T]$
    найдётся набор цен $P$, для которого пара $(h,P)$ находится в общем положении.
\end{proposition}

\begin{proof}
Множество $W_{st}$ наборов, для которых выполнен пункт~1 определения~\ref{def:general-position} для пары $(s,t)$, открыто и плотно по лемме~\ref{thm:finite}: для пары $(s,t)$ пункт~1 означает $p_t/p_s\in G_h$, а отображение $P\mapsto p_t/p_s$ непрерывно и открыто; заменив $p_t$ на $p_s\circ r$ с $r\in G_h$, близким к $p_t/p_s$, получим сколь угодно близкий набор из $W_{st}$. Поэтому $V:=U\cap\bigcap_{s\ne t}W_{st}$ --- непустое открытое множество.

Возьмём $P\in V$ и при $k=1,\dots,T$ будем по очереди заменять $p_k$ на $\lambda_kp_k$ с $\lambda_k$, близким к $1$, не выходя из открытого множества $V$; по однородности кривая $k$ при этом растягивается из начала координат в $1/\lambda_k$ раз, а остальные кривые не меняются. Покажем по индукции, что после $k$-го шага пункт~2 выполнен для всех троек с индексами $\le k$. Тройки с индексами $<k$ шаг $k$ не затрагивает. Для тройки $(r,s,k)$, $r,s<k$, множество $I_{rs}$ конечно по пункту~1 и от $\lambda_k$ не зависит, а кривая $k$ проходит через точку $x\in I_{rs}$ лишь при $\lambda_k=1/h(p_k\circ x)$. Значит, плохих $\lambda_k$ конечное число, и после $T$-го шага получаем набор из $V\subset U$ в общем положении.
\end{proof}

\begin{remark}
    Утверждения~\ref{prop:gp-stable} и~\ref{prop:gp-dense} показывают корректность
    определения~\ref{def:general-position}: множество наборов цен, находящихся с $h$ в общем
    положении, открыто и плотно в $\Rpp[2\times T]$. Иными словами, общее положение не нарушается
    при малом изменении цен и достигается сколь угодно малым изменением любого набора.
\end{remark}

\section{Разрешимость в общем положении}

\begin{definition}[Общее положение цен]\label{def:prices-gp}
    Будем говорить, что набор цен $P:=\set{p_t\in\Rpp}_{t\in[T]}$ находится в
    \textit{общем положении}, если при $s\neq t$ верно, что $p^{1}_{s}\neq p^{1}_{t}$ и $p^{2}_{s}\neq p^{2}_{t}$.
\end{definition}

\begin{proposition}\label{prop:prices-gp}
    Набор цен $P$ находится в общем положении тогда и только тогда, когда существует $h\in\hat\Phi_{2}$, для которой пара $(h,P)$ находится в общем положении.
\end{proposition}

\begin{proof}
\emph{Необходимость.} Пусть пара $(h,P)$ в общем положении. Если $p^{1}_{s}=p^{1}_{t}$ при $s\ne t$, то точка $(1/(p^{1}_{s}h(1,0)),0)$ лежит на кривых $h(p_s\circ x)=1$ и $h(p_t\circ x)=1$ и не лежит в $\Rpp$, а значит, не является трансверсальной, что противоречит пункту~1 определения~\ref{def:general-position}; так же для вторых координат.

\emph{Достаточность.} Возьмём $g(x)=\min_j\ip{\xi_j}{x}\in\hat\Phi_2$, где $\ip{\xi_j}{x}=1$ --- прямые через соседние узлы мелкого разбиения графика строго выпуклой убывающей функции. Линия $g(x)=1$ --- вписанная ломаная, и при достаточно мелком разбиении вдоль каждого её звена одна из координат меняется меньше, чем различаются одноимённые координаты любых двух цен. Тогда в общей точке $x$ кривых $s\ne t$ они проходят по разным звеньям: точки $p_s\circ x\neq p_t\circ x$ одного звена различались бы по обеим координатам. Поэтому каждое вырождение (общая вершина, касание, тройная точка) задаётся уравнением хотя бы на две разные прямые, и сдвиг одной из них $\xi_m\mapsto\xi_m+\lambda e$ превращает его в многочлен степени не выше двух от $\lambda$. Этот многочлен не равен нулю тождественно: кривая $t$ составлена из отрезков прямых $\ip{\xi_j\circ p_t}{x}=1$, и, например, для тройной точки кривых $r,s,t$ на прямых $j,k,l$ определитель $\det(\xi_j\circ p_r-\xi_k\circ p_s,\ \xi_j\circ p_r-\xi_l\circ p_t)$ при $\xi_j=(1,0)$, $\xi_k=(0,1)$, $\xi_l=0$ равен $p^1_rp^2_s\ne0$. Следовательно, меняя векторы $\xi_m$ по очереди на сколь угодно малые величины, при которых звенья остаются короткими, получаем функцию $g$, для которой пара $(g,P)$ находится в общем положении.
\end{proof}

\begin{remark}
    Наборы цен в общем положении образуют открытое всюду плотное подмножество $\Rpp[2\times T]$ из $(T!)^2$ выпуклых камер, отвечающих порядкам первых и вторых координат цен.
\end{remark}

\begin{definition}[Регулярная разрешимость]\label{def:regular-solvability}
    Пусть $P:=\set{p_t\in\Rpp}_{t\in[T]}$ --- набор цен. Будем говорить, что вектор выпусков
    $y\in\Rpp[T]$ \textit{регулярно разрешим} при ценах $P$, если существуют функция
    $h\in\hat\Phi_2$, для которой пара $(h,P)$ находится в общем положении, и полная
    абсолютно непрерывная мера $\mu$ на $\Rp$, для которой $\mu(A_t)=y_t$ при всех
    $t\in[T]$, где
    \[A_t:=\set{x\in\Rpp\mid h(p_t\circ x)<1}.\]
\end{definition}

\begin{lemma}\label{lem:signature-points}
    Рассмотрим неоклассическую функцию $h\in\Phi_{2}$ и набор цен $P:=\set{p_t\in\Rpp}_{t\in[T]}$.
    Введём \textit{спектр} пары $\mathrm{Sp}(h,P) := \left\{S\subset[T] \mid \operatorname{int}\set{x\in\Rpp\mid h(p_t\circ x)<1\Leftrightarrow t\in S} \neq \emptyset\right\}$. Тогда вектор выпусков $y\in\Rpp[T]$ принадлежит конусу $\operatorname{cone}\set{\mathbf 1_S\mid S\in\mathrm{Sp}(h,P)}$, если и только если существует полная абсолютно непрерывная мера $\mu$ на $\Rp$, для которой при всех $t\in[T]$ выполнено интегральное условие $\mu(A_t)=y_t$, где $A_t:=\set{x\in\Rpp\mid h(p_t\circ x)<1}$.
\end{lemma}

\begin{proof}
Множества $R_S:=\set{x\in\Rpp\mid h(p_t\circ x)<1\Leftrightarrow t\in S}$ разбивают $\Rpp$, и $A_t=\bigsqcup\limits_{S\ni t} R_S$, поэтому условия $\mu(A_t)=y_t$ означают $y=\sum\mu(R_S)\mathbf 1_S$. Кривые $h(p_t\circ x)=1$ имеют нулевую меру Лебега, поскольку их гомотетичные копии $h(p_t\circ x)=c$, $c>0$, попарно не пересекаются, а $R_S$ отличается от открытого множества лишь точками кривых; поэтому $\mu(R_S)=0$ при $S\notin\mathrm{Sp}(h,P)$. Обратно, разложение $y=\sum m_S \mathbf 1_S$ задаёт массу $m_S\ge0$, которую можно абсолютно непрерывно распределить по произвольному шару в $\operatorname{int}R_S$.
\end{proof}

\begin{lemma}\label{lem:gp-approx}
    Пусть набор цен $P$ находится в общем положении и $h\in\Phi_2$. Тогда существует
    $g\in\hat\Phi_2$, для которой пара $(g,P)$ находится в общем положении и
    $\mathrm{Sp}(g,P)\supseteq\mathrm{Sp}(h,P)$.
\end{lemma}

\begin{proof}
Выберем для каждого $S\in\mathrm{Sp}(h,P)$ точку $x_S\in\operatorname{int} R_S$, и введем множество $Z:=\set{p_t\circ x_S \mid t\in[T],S\in\mathrm{Sp}(h,P)}$ получив \textit{зазор} $\eta:=\min_{z\in Z}\abs{h(z)-1}>0$. Функция $H(x):=h(x)+\varepsilon\bigl(x^1+x^2+\tfrac{x^1x^2}{x^1+x^2}\bigr)$ при достаточно малом $\varepsilon>0$ отличается от $h$ на $Z$ меньше чем на $\eta/2$, а значит для каждого $z\in Z$ верно что $H(z)\gtrless1\Leftrightarrow h(z)\gtrless1$. Более того, $H\in\hat\Phi_2$ и строго вогнута на отрезке $x^1+x^2=1$, так что её линия уровня $H(x)=1$ --- график строго выпуклой убывающей функции. Выберем на этой линии точки от одной оси до другой, соединим соседние отрезками и положим $g(x):=\min_j\ip{\xi_j}{x}$, где $\ip{\xi_j}{x}=1$ --- прямые, содержащие эти отрезки. Тогда $g\in\hat\Phi_2$, линия $g(x)=1$ --- полученная ломаная, и если точки расположены достаточно часто, то $\abs{g-H}<\eta/2$ на $Z$. Сдвигая векторы $\xi_m$ на сколь угодно малые величины, приведём пару $(g,P)$ в общее положение, сохранив это неравенство. Тогда $\abs{g-h}<\eta$ на $Z$, поэтому $g(p_t\circ x_S)-1$ и $h(p_t\circ x_S)-1$ одного знака, и по непрерывности $g$ точка $x_S$ имеет спектр $S$ для $g$, то есть $S\in\mathrm{Sp}(g,P)$.
\end{proof}

\begin{theorem}\label{thm:regular-solvability}
    Пусть набор цен $P$ находится в общем положении. Тогда вектор выпусков $y\in\Rpp[T]$
    неоклассически разрешим при ценах $P$ тогда и только тогда, когда он регулярно разрешим
    при этих ценах.
\end{theorem}

\begin{proof}
По лемме~\ref{lem:signature-points} в обоих случаях речь идёт о включении $y$ в конус, натянутый на булевы векторы спектра. Пусть $y$ регулярно разрешим с функцией $h\in\hat\Phi_2\subset\Phi_2$. Во внутренности каждой области $R_S$, $S\in\mathrm{Sp}(h,P)$, есть точка вне кривых, а строгие неравенства в ней сохраняются при малом изменении цен; поэтому $\mathrm{Sp}(h,Q)\supseteq\mathrm{Sp}(h,P)$ для всех $Q$ из некоторой окрестности $P$, и $y$ неоклассически разрешим. Обратно, если $y$ неоклассически разрешим, возьмём $h\in\Phi_2$ и меру при самих ценах $P$; функция $g$ из леммы~\ref{lem:gp-approx} находится в общем положении с $P$ и $\mathrm{Sp}(g,P)\supseteq\mathrm{Sp}(h,P)$, так что $y$ регулярно разрешим.
\end{proof}

\section{Алгоритм регулярной разрешимости}

\begin{definition} [Достижимые спектры] Пусть $P:=\{p_{t}\in\Rpp\}_{t\in[T]}$ - набор цен. Будем говорить, что $t$ \textit{покрывает} пару $s, r$ если $p_{t}\in conv\set{p_{s}, p_{r}}+\Rp$. Будем говорить, что $S\subset[T]$ \textit{достижимый спектр} при ценах $P$, если ни один индекс $t\in S$ не покрывает пару индексов $s,r\notin S$. Будем обозначать как $\mathrm{Sp}(P)$ семейство таких спектров.
\end{definition}

\begin{lemma}
    Пусть $h\in\Phi_{2}$ и $P:=\{p_{t}\in\Rpp\}_{t\in[T]}$~--- набор цен. Тогда $\mathrm{Sp}(h,P)\subseteq\mathrm{Sp}(P)$.
\end{lemma}

\begin{proof}
    Предположим противное: пусть $S\in\mathrm{Sp}(h,P)$ и $t\in S$ покрывает пару $s,r\notin S$. Но тогда $A_{t}\subset A_{s}\cup A_{r}$, а значит $s$ или $r$ лежит в $S$, что противоречит условию.
\end{proof}

\begin{lemma} Пусть $P:=\{p_{t}\in\Rpp\}_{t\in[T]}$~--- набор цен. Тогда для любого $S\in\mathrm{Sp}(P)$ существует кусочно-линейная $h(x):=\min\limits_{k\in K}\langle\xi_{k}, x\rangle\in\hat\Phi_{2}$ для которой выполнено что $S\in\mathrm{Sp}(h,P)$.
\end{lemma}

\begin{proof}
    Зафиксируем $S\in\mathrm{Sp}(P)$ и построим для каждого $t\in S$ вектор $\xi_t\in\Rpp$ что $\ip{\xi_t}{p_t}<1$ и $\ip{\xi_t}{p_s}>1$ при всех $s\notin S$. Тогда $h(x):=\min\limits_{t\in S}\ip{\xi_t}{x}$ лежит в $\hat\Phi_2$, $h(p_t)\le\ip{\xi_t}{p_t}<1$ при $t\in S$ и $h(p_s)=\min\limits_{t\in S}\ip{\xi_t}{p_s}>1$ при $s\notin S$. Неравенства строгие, поэтому по непрерывности точка $x=(1,1)$ лежит в камере со спектром $S\in\mathrm{Sp}(h,P)$.

    Приведём пример таких $\xi_t$. Зафиксируем $t\in S$. Для $s\notin S$ индекс $t$ не покрывает пару $s,s$, то есть $p_s\not\le p_t$: выполнено $p^1_s>p^1_t$ или $p^2_s>p^2_t$. Пусть $\sigma_t>0$ таково, что $\sigma_t\,(p^1_s-p^1_t)+(p^2_s-p^2_t)>0$ для всех $s\notin S$. Тогда $\ip{(\sigma_t,1)}{p_t}<\min_{s\notin S}\ip{(\sigma_t,1)}{p_s}$, и вектор
    \[
        \xi_t:=\frac{2\,(\sigma_t,1)}{\ip{(\sigma_t,1)}{p_t}+\min_{s\notin S}\ip{(\sigma_t,1)}{p_s}}
    \]
    удовлетворяет $\ip{\xi_t}{p_t}<1<\ip{\xi_t}{p_s}$ при всех $s\notin S$.

    Убедимся, что такое $\sigma_t$ существует. Неравенство $\sigma_t\,(p^1_s-p^1_t)+(p^2_s-p^2_t)>0$ при $p^1_s=p^1_t$ выполнено всегда (тогда $p^2_s>p^2_t$), при $p^1_s>p^1_t$ означает $\sigma_t>\frac{p^2_t-p^2_s}{p^1_s-p^1_t}$, а при $p^1_s<p^1_t$~--- $\sigma_t<\frac{p^2_s-p^2_t}{p^1_t-p^1_s}$. Поэтому достаточно проверить, что для любых $r,s\notin S$ с $p^1_r>p^1_t>p^1_s$
    \[
        \frac{p^2_t-p^2_r}{p^1_r-p^1_t}<\frac{p^2_s-p^2_t}{p^1_t-p^1_s},
    \]
    и взять $\sigma_t$ больше нуля и всех левых частей, но меньше всех правых. Точка $q:=\lambda p_s+(1-\lambda)p_r$ с $\lambda:=\frac{p^1_r-p^1_t}{p^1_r-p^1_s}\in(0,1)$ лежит на отрезке $[p_s,p_r]$ и $q^1=p^1_t$; так как $t$ не покрывает пару $s,r$, имеем $q\not\le p_t$, то есть $q^2>p^2_t$. Подставляя $q^2=\lambda p^2_s+(1-\lambda)p^2_r$, получаем $\lambda(p^2_s-p^2_t)>(1-\lambda)(p^2_t-p^2_r)$, то есть $(p^1_r-p^1_t)(p^2_s-p^2_t)>(p^1_t-p^1_s)(p^2_t-p^2_r)$, и после деления на $(p^1_r-p^1_t)(p^1_t-p^1_s)>0$ это и есть проверяемое неравенство.
\end{proof}

\begin{lemma}[склейка]\label{lem:patch}
    Пусть $P:=\{p_{t}\in\Rpp\}_{t\in[T]}$~--- набор цен и $S_1,\dots,S_k\in\mathrm{Sp}(P)$. Тогда существует кусочно-линейная $h\in\hat\Phi_2$, для которой $S_1,\dots,S_k\in\mathrm{Sp}(h,P)$.
\end{lemma}

\begin{proof}
    Поскольку $\varnothing\in\mathrm{Sp}(h,P)$ для любой $h\in\hat\Phi_2$, пустые $S_i$ можно отбросить; считаем все $S_i$ непустыми. Для каждой пары $(i,t)$ с $t\in S_i$ возьмём по доказательству предыдущей леммы вектор $\xi_{i,t}\in\Rpp$ с $\ip{\xi_{i,t}}{p_t}<1$ и $\ip{\xi_{i,t}}{p_s}>1$ при $s\notin S_i$. Выберем $M>0$ так, чтобы $M\min(\xi^1_{i,t}p^1_s,\ \xi^2_{i,t}p^2_s)>1$ для всех таких пар $(i,t)$ и всех $s\in[T]$, и числа $0<\alpha_1<\dots<\alpha_k$ с $\alpha_{i+1}\ge M\alpha_i$. Положим $x_i:=(\alpha_i,\alpha_i^{-1})$ и $\tilde\xi_{i,t}:=(\xi^1_{i,t}\alpha_i^{-1},\ \xi^2_{i,t}\alpha_i)$. Тогда $\ip{\tilde\xi_{i,t}}{p\circ x_i}=\ip{\xi_{i,t}}{p}$ для любой цены $p$, а при $l\ne i$
    \[
        \ip{\tilde\xi_{i,t}}{p\circ x_l}=\xi^1_{i,t}p^1\,\frac{\alpha_l}{\alpha_i}+\xi^2_{i,t}p^2\,\frac{\alpha_i}{\alpha_l}\ \ge\ M\min(\xi^1_{i,t}p^1,\ \xi^2_{i,t}p^2)>1,
    \]
    поскольку одно из отношений $\alpha_l/\alpha_i$, $\alpha_i/\alpha_l$ не меньше $M$, а оба слагаемых положительны. Положим
    \[
        h(x):=\min_{i\in[k],\,t\in S_i}\ip{\tilde\xi_{i,t}}{x}\in\hat\Phi_2.
    \]
    Для $t\in S_i$ имеем $h(p_t\circ x_i)\le\ip{\tilde\xi_{i,t}}{p_t\circ x_i}=\ip{\xi_{i,t}}{p_t}<1$. Для $s\notin S_i$ каждый из векторов даёт в точке $p_s\circ x_i$ значение больше единицы: векторы $\tilde\xi_{i,t}$ той же копии~--- по выбору $\xi_{i,t}$, векторы других копий~--- по выбору $M$; значит $h(p_s\circ x_i)>1$. Все неравенства строгие, поэтому точка $x_i$ принадлежит камере со спектром $S_i\in\mathrm{Sp}(h,P)$.
\end{proof}

\begin{theorem}\label{thm:admissible-cone}
    Пусть набор цен $P$ находится в общем положении. Тогда вектор выпусков $y\in\Rpp[T]$ регулярно (равносильно, неоклассически) разрешим при ценах $P$ тогда и только тогда, когда $y\in\operatorname{cone}\set{\mathbf 1_S\mid S\in\mathrm{Sp}(P)}$.
\end{theorem}

\begin{proof}
    Равносильность регулярной и неоклассической разрешимости~--- теорема~\ref{thm:regular-solvability}. Пусть $y$ регулярно разрешим с функцией $h\in\hat\Phi_2$. По лемме~\ref{lem:signature-points} $y\in\operatorname{cone}\set{\mathbf 1_S\mid S\in\mathrm{Sp}(h,P)}$, а $\mathrm{Sp}(h,P)\subseteq\mathrm{Sp}(P)$ по лемме о достижимых спектрах. Обратно, пусть $y=\sum_{i\in[k]}m_i\mathbf 1_{S_i}$ с $m_i>0$ и $S_i\in\mathrm{Sp}(P)$. По лемме~\ref{lem:patch} найдётся $h\in\hat\Phi_2$ с $S_1,\dots,S_k\in\mathrm{Sp}(h,P)$, по лемме~\ref{lem:gp-approx}~--- функция $g\in\hat\Phi_2$ в общем положении с $P$ и $\mathrm{Sp}(g,P)\supseteq\mathrm{Sp}(h,P)\ni S_1,\dots,S_k$, и по лемме~\ref{lem:signature-points} существует полная абсолютно непрерывная мера $\mu$ с $\mu(A_t)=y_t$ для $g$. Значит $y$ регулярно разрешим.
\end{proof}

\begin{remark}
    Для $P$ в общем положении существует $h\in\hat\Phi_2$, для которой пара $(h,P)$ находится в общем положении и $\mathrm{Sp}(h,P)=\mathrm{Sp}(P)$. Действительно, применим лемму~\ref{lem:patch} ко всему семейству $\mathrm{Sp}(P)$ и получим $h_0\in\hat\Phi_2$ с $\mathrm{Sp}(h_0,P)\supseteq\mathrm{Sp}(P)$; по лемме~\ref{lem:gp-approx} найдётся $h\in\hat\Phi_2$ в общем положении с $P$ и $\mathrm{Sp}(h,P)\supseteq\mathrm{Sp}(h_0,P)$, а обратное включение $\mathrm{Sp}(h,P)\subseteq\mathrm{Sp}(P)$ даёт лемма о достижимых спектрах. Тем самым каждый регулярно разрешимый вектор выпусков реализуется одной и той же функцией $h$, и от $y$ зависит только мера.
\end{remark}

\paragraph{Алгоритм.} Теорема~\ref{thm:admissible-cone} даёт конечную процедуру проверки регулярной разрешимости $y$ при ценах $P$ в общем положении.
\begin{enumerate}
    \item \emph{Построение $\mathrm{Sp}(P)$.} Для каждой тройки индексов $t,s,r$ (допускается $s=r$) проверяем, покрывает ли $t$ пару $s,r$: условие $p_t\in\operatorname{conv}\set{p_s,p_r}+\Rp$ означает $p_t\ge\lambda p_s+(1-\lambda)p_r$ при некотором $\lambda\in[0,1]$, то есть непустоту пересечения двух отрезков значений $\lambda$, по одному на каждую координату. Затем перебираем $S\subset[T]$ и оставляем те, в которых ни один $t\in S$ не покрывает пары $s,r\notin S$. Это $O(T^3)$ проверок покрытий и $2^T$ кандидатов.
    \item \emph{Проверка принадлежности конусу.} Решаем задачу линейного программирования: найти $m_S\ge0$, $S\in\mathrm{Sp}(P)$, с $\sum_Sm_S\mathbf 1_S=y$. Если она разрешима, базисное решение даёт разложение $y$ не более чем по $T$ спектрам; вместе с леммой~\ref{lem:patch} и леммой~\ref{lem:gp-approx} это явная функция $h$ и мера, реализующие $y$. Если нет, двойственная задача даёт $\lambda\in\R^T$ с $\sum_{t\in S}\lambda_t\ge0$ при всех $S\in\mathrm{Sp}(P)$ и $\ip{\lambda}{y}<0$ (лемма Фаркаша); это опровержение: для любой реализующей пары $(h,\mu)$ из $y=\sum_S\mu(R_S)\mathbf 1_S$ и $\mathrm{Sp}(h,P)\subseteq\mathrm{Sp}(P)$ следовало бы $\ip{\lambda}{y}\ge0$. Неравенство $\sum_{\lambda_t<0}\abs{\lambda_t}y_t\le\sum_{\lambda_t>0}\lambda_ty_t$ выполнено для всех разрешимых $y$ и нарушено для данного.
\end{enumerate}
\begin{example}[проверка неразрешимости]
    Вернёмся к ценам $\hat p_1=(4,1)$, $\hat p_2=(3,3)$, $\hat p_3=(1,4)$ и вектору $y=(1,5,3)$ из раздела~1. Шаг~1: единственное покрытие~--- индекс $2$ покрывает пару $1,3$, так как $(2{,}5,\,2{,}5)\le\hat p_2$; поэтому $\mathrm{Sp}(P)=2^{[3]}\setminus\set{\set2}$. Шаг~2: задача $\sum_Sm_S\mathbf 1_S=y$, $m_S\ge0$, неразрешима, и двойственная задача даёт $\lambda=(1,-1,1)$: на каждом $S\in\mathrm{Sp}(P)$ сумма $\sum_{t\in S}\lambda_t$ неотрицательна (множество $\set2$ исключено), а $\ip{\lambda}{y}=-1<0$. Опровержение читается как неравенство $y_2\le y_1+y_3$, верное для всех разрешимых векторов и нарушенное здесь: $5>4$.
\end{example}

\begin{example}
    Опровержения не сводятся к неравенствам $y_t\le y_s+y_r$, отвечающим покрытиям. При $p_1=(107,179)$, $p_2=(59,125)$, $p_3=(26,148)$, $p_4=(66,101)$ покрытия таковы: $p_1$ доминирует $p_2,p_3,p_4$, а $2$ покрывает пару $3,4$. Вектор $y=(1,2,1,1)$ удовлетворяет всем неравенствам покрытий ($y_1\le y_2,y_3,y_4$ и $y_2\le y_3+y_4$), но шаг~2 даёт $\lambda=(-1,-1,1,1)$: всякое $S\in\mathrm{Sp}(P)$, содержащее $1$, содержит $3$ и $4$, а содержащее $2$~--- хотя бы одно из них, так что $\sum_{t\in S}\lambda_t\ge0$, тогда как $\ip{\lambda}{y}=-1$. Это неравенство $y_1+y_2\le y_3+y_4$.
\end{example}
\begin{example}[проверка разрешимости]
    Вернёмся к примеру с ценами $\hat p_1=(1,9)$, $\hat p_2=(4,4)$, $\hat p_3=(9,1)$, $\hat p_4=(10,100000)$, $\hat p_5=(20000,20000)$, $\hat p_6=(100000,10)$ и вектором $y=(5,9,6,3,1,4)$ из раздела~1. Шаг~1: цены $\hat p_1,\hat p_2,\hat p_3$ попарно несравнимы и не покрывают пар, то же верно для $\hat p_4,\hat p_5,\hat p_6$, а каждая из $\hat p_4,\hat p_5,\hat p_6$ доминирует каждую из $\hat p_1,\hat p_2,\hat p_3$; других покрытий нет. Поэтому $S\in\mathrm{Sp}(P)$ тогда и только тогда, когда из $S\cap\set{4,5,6}\ne\varnothing$ следует $\set{1,2,3}\subset S$: всего $15$ достижимых спектров, восемь подмножеств $\set{1,2,3}$ и семь множеств $\set{1,2,3}\cup U$ с непустым $U\subset\set{4,5,6}$. Шаг~2: задача разрешима,
    \[
        y=4\cdot\mathbf 1_{\set2}+1\cdot\mathbf 1_{\set3}+1\cdot\mathbf 1_{\set{1,2,3,4,5}}+2\cdot\mathbf 1_{\set{1,2,3,6}}+2\cdot\mathbf 1_{\set{1,2,3,4,6}},
    \]
    так что $y$ регулярно разрешим, что согласуется с построенной в разделе~1 неоклассической функцией. Отметим, что ни у одной CES-функции все пять спектров этого разложения не встречаются одновременно (иначе $y$ был бы CES-разрешим), тогда как по лемме~\ref{lem:patch} они реализуются одной кусочно-линейной $h$.
\end{example}
\begin{example}
    Опровержение может не иметь коэффициентов из $\set{-1,0,1}$. Пусть $p_1=(29,799)$, $p_2=(68,763)$, $p_3=(116,484)$, $p_4=(543,388)$, $p_5=(781,787)$, $p_6=(820,172)$; покрытия: $2$ покрывает пару $1,3$; $4$~--- пару $3,6$; $5$~--- пару $1,6$ и доминирует $p_2,p_3,p_4$. Вектор $y=(3,6,3,5,2,3)$ не разрешим, и $\lambda=(1,-1,2,-1,-1,1)$~--- опровержение: $y_2+y_4+y_5\le y_1+2y_3+y_6$ для всех разрешимых $y$, а здесь $13>12$. Перебор всех $3^6$ векторов $\lambda\in\set{-1,0,1}^6$ показывает, что опровержения с такими коэффициентами для этого $y$ не существует: неравенство выше задаёт грань конуса $\operatorname{cone}\set{\mathbf 1_S\mid S\in\mathrm{Sp}(P)}$, и $y$ нарушает только её.
\end{example}

\appendix
\section{Доказательство утверждения~\ref{prop:transversal}}\label{app:transversal}

Линия уровня строго положительной неоклассической функции --- график выпуклой функции $\varphi$, а супердифференциал $\partial h(x_0)$ во внутренней точке линии уровня отождествляется с отрезком опорных наклонов $\partial\varphi(t_0)$ (п.~\ref{app:A1}--\ref{app:A2}). Предельный нормальный конус графика выпуклой функции во внутренней точке вычисляется явно: это клин над отрезком наклонов вместе с двумя прямыми через его концы (п.~\ref{app:A3}). Два конуса такого вида трансверсальны ровно тогда, когда отрезки наклонов не пересекаются (п.~\ref{app:A4}); отсюда утверждение для внутренних точек. В граничной точке предельные конусы обеих кривых содержат полуплоскости, а две полуплоскости всегда имеют общий ненулевой вектор (п.~\ref{app:A5}).

\subsection{Линия уровня как график}\label{app:A1}

\begin{lemma}\label{lem:A-graph}
Пусть $h\in\hat\Phi_2$ и $b:=1/h(1,0)$. Существует непрерывная выпуклая строго убывающая функция $\varphi\colon[0,b]\to\R_+$, для которой $\varphi(0)=1/h(0,1)$, $\varphi(b)=0$ и линия уровня $h(x)=1$ в $\Rp$ совпадает с её графиком; при этом для $x^1\in[0,b]$ под графиком $h(x)<1$, а над ним $h(x)>1$, при $x^1>b$ всегда $h(x)>1$. Точка $x_0$ линии уровня лежит в $\Rpp$ тогда и только тогда, когда $x_0=(t_0,\varphi(t_0))$ с $0<t_0<b$. Функцию $\varphi$ будем называть \emph{профилем}~$h$.
\end{lemma}

\begin{proof}
Заметим сначала, что $h$ супераддитивна и строго монотонна: по однородности и вогнутости $h(x+d)=2h\bigl(\tfrac{x+d}{2}\bigr)\ge h(x)+h(d)$ при $x,d\in\Rp$, а при $d\ne0$ ещё и $h(d)>0$. При $t\in[0,b]$ функция $s\mapsto h(t,s)$ непрерывна, строго возрастает, $h(t,0)=t/b\le1$ и $h(t,s)\to\infty$ при $s\to\infty$ в силу $h(t,s)\ge s\,h(0,1)$. Значит, существует ровно одно $s\ge0$ с $h(t,s)=1$; обозначим его $\varphi(t)$. Положение точки относительно графика определяется знаком $h-1$. При $x^1>b$ имеем $h(x)\ge h(x^1,0)=x^1/b>1$. Значения $\varphi(0)=1/h(0,1)$ и $\varphi(b)=0$ получаются из однородности $h$. Строгое убывание $\varphi$ следует из строгой монотонности $h$ по первой переменной, выпуклость --- из вогнутости $h$: выпуклая комбинация двух точек графика лежит в множестве $\set{h\ge1}$, то есть не ниже графика. Непрерывность: при $s\ge0$ множества $\set{t\mid\varphi(t)>s}=\set{t\mid h(t,s)<1}$ и $\set{t\mid\varphi(t)<s}=\set{t\mid h(t,s)>1}$ открыты, так что $\varphi$ полунепрерывна и снизу, и сверху. Наконец, точка $(t_0,\varphi(t_0))$ лежит в $\Rpp$ ровно при $t_0>0$ и $\varphi(t_0)>0$, а последнее по строгому убыванию $\varphi$ и $\varphi(b)=0$ равносильно $t_0<b$.
\end{proof}

\subsection{Супердифференциал и опорные наклоны}\label{app:A2}

Напомним, что $\partial h(x_0)$ --- множество $\xi\in\R^2$ с $h(x)\le h(x_0)+\ip{\xi}{x-x_0}$ при всех $x\in\Rp$, а $\partial\varphi(t_0)$ --- множество наклонов $s$ с $\varphi(t)\ge\varphi(t_0)+s(t-t_0)$ при всех $t\in[0,b]$; для выпуклой $\varphi$ и внутренней точки $t_0$ это отрезок $[\varphi'_-(t_0),\varphi'_+(t_0)]$.

\begin{lemma}[супердифференциал через наклоны]\label{lem:A-ident}
Пусть $h\in\hat\Phi_2$, $x_0=(t_0,\varphi(t_0))$ --- внутренняя точка линии уровня $h(x)=1$. Отображение
\[
s\longmapsto\lambda(s)\,(-s,1),\qquad \lambda(s):=\bigl(\varphi(t_0)-s\,t_0\bigr)^{-1}>0,
\]
является биекцией $\partial\varphi(t_0)\to\partial h(x_0)$ с обратным $\xi\mapsto-\xi_1/\xi_2$.
\end{lemma}

\begin{proof}
Пусть $\xi\in\partial h(x_0)$. По строгому возрастанию $h$ по второй переменной $\xi_2>0$. Подстановка в опорное неравенство точек графика $x=(t,\varphi(t))$, где $h(x)=1$, даёт
\[
0\le\xi_1(t-t_0)+\xi_2(\varphi(t)-\varphi(t_0)),
\]
то есть $s:=-\xi_1/\xi_2\in\partial\varphi(t_0)$. По тождеству Эйлера $\ip{\xi}{x_0}=1$, откуда $\xi_2=\lambda(s)$ и $\xi=\lambda(s)(-s,1)$. Обратно, пусть $s\in\partial\varphi(t_0)$ и $\xi:=\lambda(s)(-s,1)$. Опорное неравенство в точке $t=b$, где $\varphi(b)=0$, даёт $s\le0$; так как $t_0>0$, имеем $\varphi(t_0)-s\,t_0\ge\varphi(t_0)>0$, то есть $\lambda(s)>0$. Далее, $\ip{\xi}{x_0}=1$, и остаётся проверить $h(x)\le\ip{\xi}{x}$ на $\Rp$. При $x=0$ это очевидно; при $x\ne0$ точка $x/h(x)$ лежит на линии уровня, и опорное неравенство для $\varphi$ в форме $\ip{\xi}{(t,\varphi(t))}\ge1$ даёт $\ip{\xi}{x}/h(x)\ge1$. Наконец, $\ip{\xi}{x-x_0}=\ip{\xi}{x}-1\ge h(x)-h(x_0)$.
\end{proof}

Отображение из леммы~\ref{lem:A-ident} зависит только от точки $x_0$: если $g\in\hat\Phi_2$ --- другая функция с $g(x_0)=1$ и профилем $\psi$, то $\psi(t_0)=\varphi(t_0)$, и функция $\lambda$ для $g$ та же. Поэтому
\begin{equation}\label{eq:A-superdiff}
\partial h(x_0)\cap\partial g(x_0)=\varnothing\iff\partial\varphi(t_0)\cap\partial\psi(t_0)=\varnothing.
\end{equation}

\subsection{Предельный нормальный конус графика}\label{app:A3}

Пусть $\varphi$ выпукла на $[a,b]$ и $t_0\in(a,b)$. Обозначим через $W$ клин, а через $\ell_\pm$ --- две прямые:
\[
W:=\set{\lambda(s,-1)\mid\lambda\ge0,\ s\in\partial\varphi(t_0)},\qquad
\ell_\pm:=\set{\mu(\varphi'_\pm(t_0),-1)\mid\mu\in\R}.
\]

\begin{lemma}\label{lem:A-cone}
Предельный нормальный конус графика $\varphi$ в точке $x_0=(t_0,\varphi(t_0))$ равен
\begin{equation}\label{eq:cone}
N(x_0)=W\cup\ell_-\cup\ell_+.
\end{equation}
\end{lemma}

\begin{proof}
\emph{Нормали Фреше.} Пусть $v$ --- нормаль Фреше к графику в точке $x_1=(t_1,\varphi(t_1))$, $t_1\in(a,b)$. Подставляя точки графика $x=(t,\varphi(t))$ в неравенство $\ip{v}{x-x_1}\le o(\abs{x-x_1})$, деля на $t-t_1$ (здесь $\abs{x-x_1}=O(\abs{t-t_1})$, так как $\varphi$ локально липшицева) и переходя к пределу при $t\downarrow t_1$ и $t\uparrow t_1$, получаем
\begin{equation}\label{eq:frechet}
v_1+v_2\,\varphi'_+(t_1)\le0\le v_1+v_2\,\varphi'_-(t_1).
\end{equation}
С другой стороны, клин $W$ состоит из нормалей Фреше в $x_0$: для $v=\lambda(s,-1)$ неравенство $\ip{v}{x-x_0}\le0$ на графике следует из опорного неравенства $\varphi(t)\ge\varphi(t_0)+s(t-t_0)$. Если $\varphi$ дифференцируема в $t_1$, то нормалью Фреше в $x_1$ является и вектор $(-\varphi'(t_1),1)$: на графике
\[
\ip{v}{x-x_1}=\varphi(t)-\varphi(t_1)-\varphi'(t_1)(t-t_1)=o(\abs{t-t_1}).
\]

\emph{Включение $\supset$.} Клин $W$ содержится в $N(x_0)$, поскольку состоит из нормалей Фреше. Выпуклая функция дифференцируема всюду, кроме счётного множества, поэтому найдутся точки $t_n\downarrow t_0$, в которых $\varphi$ дифференцируема; поскольку правая производная выпуклой функции непрерывна справа, $\varphi'(t_n)\to\varphi'_+(t_0)$. Значит, вектор $(-\varphi'_+(t_0),1)$ является пределом нормалей Фреше $(-\varphi'(t_n),1)$, то есть предельной нормалью; вместе с $(\varphi'_+(t_0),-1)\in W$ это даёт $\ell_+\subset N(x_0)$. Для $\ell_-$ аналогично, с $t_n\uparrow t_0$ и непрерывностью $\varphi'_-$ слева.

\emph{Включение $\subset$.} Пусть $v=\lim v_n$, где $v_n$ --- нормали Фреше в точках графика с абсциссами $t_n\to t_0$. Переходя к подпоследовательности, считаем, что все $t_n>t_0$, все $t_n<t_0$ или все $t_n=t_0$. В первых двух случаях предельный переход в~\eqref{eq:frechet} даёт $v_1+v_2\varphi'_\pm(t_0)=0$, то есть $v\in\ell_\pm$. В третьем случае $v$ удовлетворяет~\eqref{eq:frechet} в точке $t_0$: при $v_2\le0$ это означает $v\in W$, а при $v_2>0$ из $\varphi'_+(t_0)\le-v_1/v_2\le\varphi'_-(t_0)$ следует $\varphi'_-(t_0)=\varphi'_+(t_0)$, и $v\in\ell_+$.
\end{proof}

\subsection{Трансверсальность двух графиков и линий уровня}\label{app:A4}

\begin{lemma}\label{lem:A-cones}
Пусть $[d_-,d_+]$ и $[e_-,e_+]$ --- два отрезка, а $W_d,\ell^d_\pm$ и $W_e,\ell^e_\pm$ --- клинья и прямые, построенные по ним так же, как $W,\ell_\pm$ по отрезку $\partial\varphi(t_0)$ в п.~\ref{app:A3}. Конусы
\[
C_d:=W_d\cup\ell^d_-\cup\ell^d_+\quad\text{и}\quad C_e:=W_e\cup\ell^e_-\cup\ell^e_+
\]
трансверсальны, то есть $C_d\cap(-C_e)=\set0$, тогда и только тогда, когда $[d_-,d_+]\cap[e_-,e_+]=\varnothing$.
\end{lemma}

\begin{proof}
Клин $-W_e$ состоит из векторов с неотрицательной второй координатой, $W_d$ --- с неположительной, поэтому ненулевой общий вектор не может лежать в обоих клиньях. Прямые симметричны относительно нуля, и сравнение координат в остальных случаях даёт: $d_\pm\in[e_-,e_+]$, либо $e_\pm\in[d_-,d_+]$, либо $d_\pm=e_\pm$ --- во всех случаях отрезки пересекаются. Обратно, если отрезки пересекаются, то конец одного лежит в другом, скажем $d_+\in[e_-,e_+]$; тогда $v:=(-d_+,1)\in\ell^d_+$ и $-v=(d_+,-1)\in W_e$.
\end{proof}

\begin{proof}[Доказательство утверждения~\ref{prop:transversal} для внутренних точек]
По лемме~\ref{lem:A-graph} линии уровня $h(x)=1$ и $g(x)=1$ совпадают с графиками профилей $\varphi$ и $\psi$, и $x_0=(t_0,\varphi(t_0))=(t_0,\psi(t_0))$ --- внутренняя точка обоих графиков. По лемме~\ref{lem:A-cone} их предельные нормальные конусы в $x_0$ являются объединениями клина и двух прямых, построенных по отрезкам $\partial\varphi(t_0)$ и $\partial\psi(t_0)$. По лемме~\ref{lem:A-cones} трансверсальность равносильна $\partial\varphi(t_0)\cap\partial\psi(t_0)=\varnothing$, а по~\eqref{eq:A-superdiff} --- условию $\partial h(x_0)\cap\partial g(x_0)=\varnothing$.
\end{proof}

\subsection{Граничные точки}\label{app:A5}

\begin{lemma}\label{lem:A-endpoint}
Пусть $\varphi$ выпукла на $[a,b]$, $a<b$, и $\varphi\ge\varphi(b)$ на $[a,b]$. Тогда левая производная $\varphi'_-(b)$ существует и конечна, а полуплоскость $\set{v\mid v_1+\varphi'_-(b)\,v_2\ge0}$ содержится в конусе нормалей Фреше к графику $\varphi$ в точке $x_0:=(b,\varphi(b))$.
\end{lemma}

\begin{proof}
Наклоны хорд $(\varphi(b)-\varphi(t))/(b-t)$, $t\in(a,b)$, неположительны и не убывают по $t$ в силу выпуклости, поэтому имеют конечный предел при $t\uparrow b$, равный их супремуму; это и есть $\varphi'_-(b)$. Для $v$ из указанной полуплоскости и точки графика $x=(t,\varphi(t))$, $t<b$, имеем
\[
\ip{v}{x-x_0}=(t-b)\Bigl(v_1+v_2\,\frac{\varphi(b)-\varphi(t)}{b-t}\Bigr),
\]
а выражение в скобках стремится к $v_1+v_2\varphi'_-(b)\ge0$, так что $\ip{v}{x-x_0}\le o(\abs{x-x_0})$.
\end{proof}

\begin{proof}[Доказательство утверждения~\ref{prop:transversal} для граничных точек]
Пусть $x_0\in\partial\Rpp$ лежит на обеих линиях уровня, и $N_h(x_0)$, $N_g(x_0)$ --- предельные нормальные конусы линий уровня $h(x)=1$ и $g(x)=1$ в точке $x_0$. Рассмотрим сначала ось $x^2=0$. Линия $h(x)=1$ пересекает её только в точке $(b,0)$, $b=1/h(1,0)$, так как $h(t,0)=t\,h(1,0)$; значит, $x_0=(b,0)$ --- правый конец графика профиля $\varphi$. Поскольку $g(x_0)=1$, точка $x_0$ является правым концом и графика профиля $\psi$ функции $g$. Так как $\varphi\ge0=\varphi(b)$, по лемме~\ref{lem:A-endpoint} конус $N_h(x_0)$ содержит полуплоскость $\set{v_1+s_hv_2\ge0}$, $s_h:=\varphi'_-(b)$, а $-N_g(x_0)$ содержит полуплоскость $\set{v_1+s_gv_2\le0}$, $s_g:=\psi'_-(b)$. Две такие полуплоскости имеют общий ненулевой вектор: $(-s_h,1)$ при $s_g\le s_h$ и $(s_g,-1)$ при $s_g>s_h$. Значит, $N_h(x_0)\cap(-N_g(x_0))\ne\set0$.

Ось $x^1=0$ сводится к разобранному случаю перестановкой координат $\sigma(x^1,x^2):=(x^2,x^1)$. Это линейная изометрия, переводящая $\Rp$ в себя, поэтому она переводит конусы Фреше и предельные конусы множества $A$ в соответствующие конусы множества $\sigma(A)$, и трансверсальность пары множеств сохраняется при применении $\sigma$. Кроме того, $h\circ\sigma\in\hat\Phi_2$ вместе с $h$, а $\sigma$ переводит линию уровня $h(x)=1$ в линию уровня $h(x^2,x^1)=1$. После перестановки точка $x_0$ попадает на ось $x^2=0$.
\end{proof}

Отметим, что бесконечная производная профиля здесь не возникает: она возможна только в левом конце графика, а левый конец после перестановки координат становится правым, где производная конечна по лемме~\ref{lem:A-endpoint}.

{\raggedright\paragraph{Соответствие с формализацией.} Доказательство формализовано в файле \lean{Transversality.lean}. Конусы и трансверсальность~--- \lean{frechetNormal}, \lean{limitingNormal}, \lean{SetTransversal}; лемма~\ref{lem:A-graph}~--- \lean{levelCurve_eq_graphOn}, \lean{curveFun_convexOn}, \lean{mem_Ioo_of_mem_levelCurve} (профиль определён через \lean{profile} из \lean{Profile.lean}); лемма~\ref{lem:A-ident}~--- \lean{neg_div_mem_subdiffOn}, \lean{mem_superdiff_of_mem_subdiffOn} (тождество Эйлера и положительность $\xi_2$~--- \lean{ip_eq_of_mem_superdiff}, \lean{snd_pos_of_mem_superdiff}), а эквивалентность~\eqref{eq:A-superdiff}~--- \lean{disjoint_superdiff_iff}; лемма~\ref{lem:A-cone}~--- \lean{limitingNormal_graph_eq}, где неравенства~\eqref{eq:frechet} суть \lean{frechet_graph_bound_right}, \lean{frechet_graph_bound_left}, клин~--- \lean{wedge_subset_frechet}, нормаль в точке дифференцируемости~--- \lean{neg_grad_mem_frechet}, плотность точек дифференцируемости~--- \lean{exists_differentiableAt_Ioo}, сходимость односторонних производных~--- \lean{tendsto_derivs_nhdsGT}, \lean{tendsto_derivs_nhdsLT}; лемма~\ref{lem:A-cones}~--- \lean{transversal_cones_iff}; утверждение для внутренних точек~--- \lean{transversal_graphs_iff}, \lean{transversal_levelCurves_iff}; лемма~\ref{lem:A-endpoint}~--- \lean{exists_hasDerivWithinAt_Iio_endpoint}, \lean{halfplane_subset_frechet_endpoint}; общий вектор двух полуплоскостей~--- \lean{exists_ne_zero_of_halfplanes}; граничные точки~--- \lean{not_setTransversal_of_snd_eq_zero} и, после перестановки координат \lean{sw} с леммами об инвариантности конусов и класса $\hat\Phi_2$, \lean{not_setTransversal_of_not_mem_orthant}.\par}

\section{Доказательство леммы~\ref{thm:finite}}\label{app:finite}

Точку пересечения кривых называем трансверсальной, если она лежит в $\Rpp$ и супердифференциалы двух функций в ней не пересекаются; по утверждению~\ref{prop:transversal} это трансверсальность в смысле предельных нормальных конусов. Растяжением координат задача сводится к случаю $h(1,0)=h(0,1)=1$, $p^1>1>p^2$ (п.~\ref{app:B1}). Линия уровня такой функции --- график её профиля $\varphi$ на $[0,1]$, и точки пересечения кривых взаимно однозначно соответствуют решениям уравнения $\alpha\varphi(y)=\varphi(\beta y)$ (п.~\ref{app:B2}). Для этого уравнения трансверсальные решения изолированы, а параметры $(\alpha,\beta)$, при которых все решения трансверсальны, образуют открытое всюду плотное множество (п.~\ref{app:B3}). Наконец, трансверсальность решения уравнения равносильна трансверсальности соответствующей точки пересечения (п.~\ref{app:B4}).

\subsection{Обозначения и сведение к основному случаю}\label{app:B1}

Для $p\in\Rpp$ обозначим $D_px:=p\circ x$ и $p^{-1}:=(1/p^1,1/p^2)$; все $D_p$ --- диагональные линейные автоморфизмы $\Rp$, попарно коммутирующие, $D_{p^{-1}}=D_p^{-1}$. Положим $h_p:=h\circ D_p$ и
\[
I_h(p):=\set{x\in\Rp\mid h(x)=1,\ h(p\circ x)=1}.
\]
Супердифференциал $\partial h(x_0)$ понимаем как в п.~\ref{app:A2}. Точку $x_0\in I_h(p)$ называем \emph{трансверсальной}, если $x_0\in\Rpp$ и $\partial h(x_0)\cap\partial h_p(x_0)=\varnothing$; поскольку $h_p\in\hat\Phi_2$ вместе с $h$, по утверждению~\ref{prop:transversal} это трансверсальность кривых $h(x)=1$ и $h_p(x)=1$ в смысле предельных нормальных конусов. Пользуемся также супераддитивностью и строгой монотонностью $h$ (см. доказательство леммы~\ref{lem:A-graph}).

\paragraph{Растяжения.} Линейная замена $x\mapsto D_cx$, $c\in\Rpp$, переводит $\Rp$ на себя и коммутирует с $D_p$, поэтому $I_{h\circ D_c}(p)=D_c^{-1}I_h(p)$, а супердифференциалы преобразуются по цепному правилу $\partial(g\circ D_c)(x)=D_c\,\partial g(D_cx)$. Значит, конечность $I_h(p)$ и трансверсальность его точек инвариантны относительно замены $h\mapsto h\circ D_c$. Аналогично замена $x\mapsto D_px$ меняет местами кривые $h(x)=1$ и $h_p(x)=1$: $x\in I_h(p)\iff D_px\in I_h(p^{-1})$, и трансверсальность точки $D_px$ для пары $(h,h_{p^{-1}})$ равносильна трансверсальности $x$ для пары $(h,h_p)$.

\paragraph{Нормировка.} Для $h\in\hat\Phi_2$ положим $c:=(h(1,0),h(0,1))$ и $h_0:=h\circ D_{c^{-1}}$; тогда $h_0\in\hat\Phi_2$, $h_0(1,0)=h_0(0,1)=1$ и $h=h_0\circ D_c$. Поэтому лемму~\ref{thm:finite} достаточно доказать для $h_0$, и далее считаем $h(1,0)=h(0,1)=1$.

\paragraph{Вырожденные случаи.} Если $p^1,p^2>1$, то $I_h(p)=\varnothing$: при $h(x)=1$ и $m:=\min\set{p^1,p^2}$ имеем $p\circ x\ge mx$ покоординатно, откуда по монотонности и однородности $h(p\circ x)\ge h(mx)=m>1$. Аналогично при $p^1,p^2<1$ получаем $h(p\circ x)\le\max\set{p^1,p^2}<1$. Случай $p^1<1<p^2$ сводится к случаю $p^1>1>p^2$ заменой $x\mapsto D_px$. Остаётся \emph{основной случай} $p^1>1>p^2$.

\subsection{Профиль и одномерное уравнение}\label{app:B2}

По лемме~\ref{lem:A-graph} для нормированной $h\in\hat\Phi_2$ линия уровня $h(x)=1$ совпадает с графиком профиля $\varphi\colon[0,1]\to[0,1]$, $\varphi(0)=1$, $\varphi(1)=0$; при этом для $x^1\in[0,1]$ под графиком $h(x)<1$, а над ним $h(x)>1$, при $x^1>1$ всегда $h(x)>1$.

Пусть $p^1>1>p^2$ и $\alpha:=1/p^2$, $\beta:=1/p^1$, так что $(\alpha,\beta)\in P:=(1,+\infty)\times(0,1)$. Точка $x=(t,\varphi(t))$ линии уровня лежит на второй кривой тогда и только тогда, когда $p\circ x=(p^1t,\ p^2\varphi(t))$ лежит на графике, то есть $p^1t\le1$ и $\varphi(p^1t)=p^2\varphi(t)$. В переменной $y:=p^1t\in[0,1]$ это уравнение
\begin{equation}\label{eq:B-profile}
\alpha\varphi(y)=\varphi(\beta y),\qquad y\in[0,1],
\end{equation}
и $y\mapsto(y/p^1,\varphi(y/p^1))$ --- биекция множества его решений $S(\alpha,\beta)$ на $I_h(p)$. Концы отрезка решениями не являются: $\alpha\varphi(0)=\alpha>1=\varphi(0)$ и $\alpha\varphi(1)=0<\varphi(\beta)$. Поэтому все точки $I_h(p)$ лежат в $\Rpp$. Отметим, что условие $\beta<1$ нужно, чтобы аргумент $\beta y$ оставался в области определения $\varphi$, а строгое убывание $\varphi$ существенно: если бы $\varphi$ обращалась в нуль на отрезке $[s,1]$, то при $\beta>s$ все $y\in[s/\beta,1]$ были бы решениями~\eqref{eq:B-profile} при любом $\alpha$.

\subsection{Одномерная задача}\label{app:B3}

Для $y\in(0,1)$ субдифференциал $\partial\varphi(y)=[\varphi'_-(y),\varphi'_+(y)]$ --- отрезок опорных наклонов. Напомним стандартные факты о выпуклых функциях одной переменной: $\partial\varphi(y)\subset[-1/y,0]$ (опорные неравенства в точках $0$ и $1$), отображение $y\mapsto\partial\varphi(y)$ имеет замкнутый график, а $\varphi$ дифференцируема всюду, кроме не более чем счётного множества, причём там, где она дифференцируема, $\partial\varphi(y)=\set{\varphi'(y)}$.

\begin{definition}\label{def:B-transversal}
Решение $y\in S(\alpha,\beta)$ \emph{трансверсально}, если $\alpha\,\partial\varphi(y)\cap\beta\,\partial\varphi(\beta y)=\varnothing$; параметр $(\alpha,\beta)\in P$ \emph{трансверсален}, если трансверсальны все его решения. Пару $(y,\beta)$, $y,\beta\in(0,1)$, назовём \emph{критической}, если $\beta u\,\varphi(y)=\varphi(\beta y)\,w$ для некоторых $u\in\partial\varphi(\beta y)$, $w\in\partial\varphi(y)$.
\end{definition}

Для решения $y$ уравнения~\eqref{eq:B-profile} критичность пары $(y,\beta)$ равносильна нетрансверсальности $y$: достаточно разделить на $\varphi(y)>0$ и подставить $\varphi(\beta y)=\alpha\varphi(y)$. Преимущество критичности в том, что она не зависит от $\alpha$ и задаёт замкнутое условие.

\begin{lemma}[трансверсальные решения изолированы]\label{lem:B-finite}
Если $y_0$ --- предельная точка множества решений уравнения~\eqref{eq:B-profile}, то пара $(y_0,\beta)$ критическая. В частности, у трансверсального параметра конечно много решений.
\end{lemma}

\begin{proof}
Разность $\alpha\varphi(z)-\varphi(\beta z)$ непрерывна, поэтому обращается в нуль не только в решениях $y_n\downarrow y_0$ (случай $y_n\uparrow y_0$ аналогичен), но и в $y_0$; при этом $y_0\in(0,1)$, так как $0$ и $1$ решениями не являются. Тогда её правая производная в $y_0$ равна нулю: $\alpha\varphi'_+(y_0)=\beta\varphi'_+(\beta y_0)$, а односторонние производные лежат в субдифференциалах, так что $(y_0,\beta)$ критическая. Множество решений замкнуто в $[0,1]$; будь оно бесконечным, у него нашлась бы точка накопления --- решение, которое по доказанному не трансверсально.
\end{proof}

\begin{lemma}[локализация решений]\label{lem:B-local}
Для $1<a\le b$ и $0<\beta_1\le\beta_2<1$ найдётся отрезок $[c,d]\subset(0,1)$, содержащий $S(\alpha,\beta)$ при всех $\alpha\in[a,b]$, $\beta\in[\beta_1,\beta_2]$.
\end{lemma}

\begin{proof}
По непрерывности $\varphi$ в нуле при малых $y$ выполнено $\alpha\varphi(y)\ge a\varphi(y)>1\ge\varphi(\beta y)$, а по непрерывности в единице при $y$, близких к $1$, --- $\alpha\varphi(y)\le b\varphi(y)<\varphi(\beta_2)\le\varphi(\beta y)$, поскольку $\beta y\le\beta_2$ и $\varphi$ убывает.
\end{proof}

\begin{lemma}[плохие параметры содержатся в компакте]\label{lem:B-compact}
Пусть $K:=[c,d]\times[\beta_1,\beta_2]$, где $[c,d]$ взят из леммы~\ref{lem:B-local}, $Z\subset K$ --- множество критических пар и $F(y,\beta):=\varphi(\beta y)/\varphi(y)$. Тогда $Z$ компактно, $F$ непрерывна на $K$, множество
\[
B:=\set{(F(y,\beta),\beta)\mid(y,\beta)\in Z}
\]
компактно, и параметр $(\alpha,\beta)\in[a,b]\times[\beta_1,\beta_2]$ трансверсален тогда и только тогда, когда он не лежит в $B$.
\end{lemma}

\begin{proof}
Субградиенты $u\in\partial\varphi(\beta y)$, $w\in\partial\varphi(y)$ для $(y,\beta)\in K$ равномерно ограничены, а график субдифференциала замкнут, поэтому условие критичности сохраняется при предельном переходе, и $Z$ замкнуто в компакте $K$. Непрерывность $F$ следует из $\varphi\ge\varphi(d)>0$ на $[c,d]$, а $B$ --- образ компакта при непрерывном отображении. Если решение $y$ параметра $(\alpha,\beta)$ не трансверсально, то $y\in[c,d]$ по лемме~\ref{lem:B-local}, пара $(y,\beta)$ критическая и $F(y,\beta)=\alpha$, то есть $(\alpha,\beta)\in B$. Обратно, если $(\alpha,\beta)=(F(y,\beta),\beta)$ для критической пары $(y,\beta)\in Z$, то $\alpha\varphi(y)=\varphi(\beta y)$, и $y$ --- нетрансверсальное решение.
\end{proof}

\begin{lemma}[мера критических значений на срезе]\label{lem:B-sard}
При фиксированном $\beta_0$ множество $\set{F(y,\beta_0)\mid (y,\beta_0)\in Z}$ имеет меру нуль.
\end{lemma}

\begin{proof}
Пусть $F_0:=F(\cdot,\beta_0)$. Критические $y$, для которых $\varphi$ не дифференцируема в $y$ или в $\beta_0y$, образуют счётное множество, и его образ счётен. В остальных критических точках субдифференциалы одноточечны, условие критичности означает $\beta_0\varphi'(\beta_0y)\varphi(y)=\varphi(\beta_0y)\varphi'(y)$, а это в точности равенство $F_0'(y)=0$ по правилу дифференцирования частного. Образ множества точек с нулевой производной имеет меру нуль по лемме Сарда~\cite{Sard1942}.
\end{proof}

\begin{lemma}[общее положение для одномерной задачи]\label{lem:B-1d}
Трансверсальные параметры образуют открытое всюду плотное подмножество $P$.
\end{lemma}

\begin{proof}
Возьмём $(\alpha_0,\beta_0)\in P$ и прямоугольник $R:=[\alpha_0-r,\alpha_0+r]\times[\beta_1,\beta_2]\subset P$ с $\beta_1<\beta_0<\beta_2$, сколь угодно малый; применим к нему леммы~\ref{lem:B-local} и~\ref{lem:B-compact}. Множество $B$ замкнуто, поэтому $\operatorname{int}R\setminus B$ открыто и по лемме~\ref{lem:B-compact} состоит в точности из трансверсальных параметров прямоугольника. Если $(\alpha_0,\beta_0)$ трансверсален, то он не лежит в $B$, и $\operatorname{int}R\setminus B$ --- его окрестность из трансверсальных параметров; это даёт открытость. Для плотности воспользуемся леммой~\ref{lem:B-sard}: множество $E:=\set{F(y,\beta_0)\mid(y,\beta_0)\in Z}$ имеет меру нуль и не покрывает интервал $(\alpha_0-r,\alpha_0+r)$; для $\alpha_1\in(\alpha_0-r,\alpha_0+r)\setminus E$ имеем $(\alpha_1,\beta_0)\in\operatorname{int}R\setminus B$.
\end{proof}

\subsection{Перевод трансверсальности}\label{app:B4}

\begin{lemma}\label{lem:B-translate}
Пусть $h(1,0)=h(0,1)=1$, $p^1>1>p^2$, $(\alpha,\beta)=(1/p^2,1/p^1)$ и $y_0$ --- решение уравнения~\eqref{eq:B-profile}. Точка $x_0:=(t_0,\varphi(t_0))\in I_h(p)$, где $t_0:=\beta y_0$, трансверсальна тогда и только тогда, когда трансверсально решение $y_0$.
\end{lemma}

\begin{proof}
По п.~\ref{app:B2} $x_0\in I_h(p)\cap\Rpp$. Функция $h_p$ принадлежит $\hat\Phi_2$, $h_p(1,0)=p^1$, и по лемме~\ref{lem:A-graph} её линия уровня совпадает с графиком функции $\psi(z):=\varphi(p^1z)/p^2$ на $[0,1/p^1]$: действительно, $h_p(z,\psi(z))=h(p^1z,\varphi(p^1z))=1$. Точка $x_0$ --- внутренняя точка обоих графиков, поэтому по~\eqref{eq:A-superdiff} условие $\partial h(x_0)\cap\partial h_p(x_0)=\varnothing$ равносильно $\partial\varphi(t_0)\cap\partial\psi(t_0)=\varnothing$. Наконец, $\partial\psi(t_0)=(p^1/p^2)\,\partial\varphi(y_0)$ по правилу дифференцирования при линейной замене переменной, и после деления на $p^1$ условие принимает вид $\beta\,\partial\varphi(\beta y_0)\cap\alpha\,\partial\varphi(y_0)=\varnothing$.
\end{proof}

\subsection{Доказательство леммы~\ref{thm:finite}}\label{app:B5}

Прямые $p^1=1$ и $p^2=1$ не пересекаются с $G_h$: при $p^1=1$ точка $(1/h(1,0),0)$ лежит на обеих кривых, но не в $\Rpp$, и потому не трансверсальна; так же при $p^2=1$ с точкой $(0,1/h(0,1))$. Камеры $\set{p^1,p^2>1}$ и $\set{p^1,p^2<1}$ целиком лежат в $G_h$: по п.~\ref{app:B1} пересечений нет вовсе. Пусть $Q:=\set{p^1>1>p^2}$; отображение $\tau(p):=(1/p^2,1/p^1)$ --- гомеоморфизм $Q$ на $P$. Для $p\in Q$ решения уравнения~\eqref{eq:B-profile} с параметром $\tau(p)$ взаимно однозначно соответствуют точкам $I_h(p)$ (п.~\ref{app:B2}), поэтому по леммам~\ref{lem:B-finite} и~\ref{lem:B-translate} условие $p\in G_h$ равносильно трансверсальности параметра $\tau(p)$, и по лемме~\ref{lem:B-1d} множество $G_h\cap Q$ открыто и плотно в $Q$. Замена $x\mapsto D_px$ из п.~\ref{app:B1} показывает, что $p\in G_h\iff p^{-1}\in G_h$, а инверсия $p\mapsto p^{-1}$ переводит камеру $\set{p^1<1<p^2}$ на $Q$; значит, и в этой камере $G_h$ открыто и плотно. Итак, $G_h$ --- объединение открытых плотных подмножеств четырёх открытых камер, объединение которых плотно в $\Rpp$; следовательно, $G_h$ открыто и всюду плотно. \qed

{\raggedright\paragraph{Соответствие с формализацией.} Доказательство формализовано в файлах \lean{Neoclassical.lean}, \lean{Profile.lean}, \lean{FiniteIntersections.lean}, \lean{LevelCurves.lean}; части, относящиеся к открытости, --- в \lean{GeneralPosition.lean}. Трансверсальность точки в коде~--- предикат \lean{TransversalAt}, супердифференциал~--- \lean{superdiff}; растяжения и инверсия~--- \lean{superdiff_comp_hadamard}, \lean{levelIntersections_eq_preimage}, \lean{TransversalAt.of_hinv}; вырожденные случаи~--- \lean{levelIntersections_eq_empty_of_one_lt}, \lean{levelIntersections_eq_empty_of_lt_one}; профиль и сведение к одной переменной~--- \lean{profile_isProfile}, \lean{ofProfile_profile}, \lean{levelIntersections_ofProfile}; определение~\ref{def:B-transversal}~--- \lean{Transversal}, \lean{IsCritical}; лемма~\ref{lem:B-finite}~--- \lean{isCritical_of_frequently}, \lean{finite_of_transversal}; лемма~\ref{lem:B-local}~--- \lean{exists_Icc_solutions_subset}; лемма~\ref{lem:B-compact}~--- \lean{isCompact_critical}; лемма~\ref{lem:B-sard}~--- \lean{volume_image_critical_slice} (в коде вместо счётности точек недифференцируемости использованы теорема Радемахера~\cite{Rademacher1919,EvansGariepy2015} и липшицевость $F_0$); лемма~\ref{lem:B-1d}~--- \lean{exists_open_transversal} (плотность) и \lean{IsProfile.eventually_transversal} (открытость); лемма~\ref{lem:B-translate}~--- \lean{transversalAt_ofProfile} (импликация от решения к точке, прямой подстановкой в опорные неравенства) и \lean{transversal_of_isGoodPrice} (обратная импликация, через \lean{disjoint_superdiff_iff}); множество $G_h$~--- \lean{IsGoodPrice}; лемма~\ref{thm:finite}~--- \lean{exists_open_finite_levelIntersections} (плотность) и \lean{eventually_isGoodPrice} (открытость).\par}

\begin{thebibliography}{9}
\bibitem{Kruger2006} A.~Y.~Kruger, About regularity of collections of sets,
  Set-Valued Anal. 14 (2006) 187--206.
\bibitem{Ioffe2017} A.~D.~Ioffe, Variational Analysis of Regular Mappings:
  A View from Variational Analysis, Springer, 2017.
\bibitem{Mordukhovich2006} B.~S.~Mordukhovich, Variational Analysis and Generalized
  Differentiation I: Basic Theory, Springer, 2006.
\bibitem{Sard1942} A.~Sard, The measure of the critical values of differentiable maps,
  Bull. Amer. Math. Soc. 48 (1942) 883--890.
\bibitem{Rademacher1919} H.~Rademacher, \"Uber partielle und totale Differenzierbarkeit von
  Funktionen mehrerer Variabeln und \"uber die Transformation der Doppelintegrale,
  Math. Ann. 79 (1919) 340--359.
\bibitem{EvansGariepy2015} L.~C.~Evans, R.~F.~Gariepy, Measure Theory and Fine Properties of
  Functions, revised ed., CRC Press, 2015.
\end{thebibliography}

\end{document}
